% !TeX root = ../../main.tex
\section{计算需求差异的形式化分析}
\subsection{输出要求与近似预算的层次}
对计算任务定义输出要求，有助于确定哪些近似可以接受。检索的输入为固定数据库和查询，契约包括得分函数、目标 $K$、并列值处理以及召回阈值。CG 的契约则包括矩阵类别、右端项、初始解、停止准则和失败处理。两类契约中的阈值不能互换，例如检索 Recall 为0.99 并不能对应数值求解相对残差为0.01。

近似预算可以作用于多个层次。数值预算决定每个被处理元素使用多少位；范围预算决定处理多少候选或簇；控制预算决定多久重新选择精度。预算并非越小越好，因为减少一项开销可能增加另一项。例如更低位宽需要更大的候选集合，较少策略调用可能延迟必要的精度提升。分析应先明确预算影响的对象，再测量其对完整任务的影响。

\begin{table}[htbp]\centering\caption{不同预算的作用对象与补偿路径}
\begin{tabularx}{\textwidth}{p{2.2cm}XXX}\toprule
预算&直接节省&可能引入的代价&可用补偿方式\\\midrule
数值位宽&数值字段与算术资源&得分偏差、局部扰动&候选扩展、高精度重算\\
候选规模&精排读取与计算&相关条目被截断&扩大局部候选、重新筛选\\
簇探测数&数据库扫描量&相关区域被遗漏&提高探测预算、改善聚类\\
非零保留比例&非零记录数量&向量方向及得分变化&校准保留比例、保留原始副本\\
动作复用间隔&控制调用次数&响应状态变化滞后&缩短间隔、重新画像\\\bottomrule
\end{tabularx}\end{table}

\subsection{排序扰动的充分条件与候选扩展}
设近似分数为 $\hat y_i=y_i+e_i$，且 $|e_i|\leq\varepsilon_i$。真实 Top-K 中的行 $i$ 与集合外的行 $j$ 保持正确相对次序的充分条件是
\begin{equation}
 y_i-y_j>\varepsilon_i+\varepsilon_j.
\end{equation}
原因是近似差值的下界为 $\hat y_i-\hat y_j\geq y_i-y_j-\varepsilon_i-\varepsilon_j$。若右侧为正，近似次序不变。该条件关注排序间隔而不是分数本身的相对误差，因此两个具有相同平均误差的数据集可能拥有不同的候选召回。

将预筛规模从 $K$ 扩大到 $L$，需要关心真实 Top-K 是否被推到近似第 $L$ 名之后。对任意真实 Top-K 行，若最多只有 $L-1$ 个候选可能在扰动后超过它，则该行仍会保留。实际系统通常没有全部行的可用误差界，故通过验证查询估计覆盖。这一推导揭示了候选扩展的作用机制；实际配置的候选覆盖仍通过独立查询校准与验证。

当采用分布式局部候选时，风险又增加一层。即使全局近似 Top-$L$覆盖目标，局部分区容量不足仍可能删除重要行。质量审计应区分数值扰动和容量截断，分别使用“精确分数加局部选择”和“近似分数加全局选择”的对照，避免两种损失互相遮蔽。

\subsection{稀疏性与传输效率的分离}
逻辑稀疏性用 $z/(mn)$ 描述，而物理效率还取决于填充与事务粒度。设有效非零记录字节数为 $B_{\mathrm{valid}}$，实际读取为 $B_{\mathrm{actual}}$，则有效载荷率为
\begin{equation}
 \eta_{\mathrm{payload}}=B_{\mathrm{valid}}/B_{\mathrm{actual}}.
\end{equation}
剪枝使 $B_{\mathrm{valid}}$下降，但若每行仍需读取完整固定长度块，$B_{\mathrm{actual}}$可能不同比下降。类似地，过小的簇虽减少候选，却可能增加地址表和事务启动占比。

设每条总线记录包含数值 $b_v$、索引 $b_i$ 和边界信息 $b_f$ 位，总线宽度为 $W$，单包容量是
\begin{equation}
 P=\left\lfloor W/(b_v+b_i+b_f)\right\rfloor.
\end{equation}
该表达式是理想打包上限。实际实现还需考虑控制标记、跨包延续和对齐。比较不同格式时，应同时列出理论容量、实际包格式和测得带宽，避免把位级上限作为可持续性能。

\subsection{并行任务的不均衡与关键路径}
平均工作量只描述总体资源需求，完成时间还取决于任务最大值。若每核速率为 $v$，忽略其他开销，理想均衡时间为 $\sum_pW_p/(Pv)$，实际受最重核限制为 $\max_pW_p/v$。据此可定义不均衡因子
\begin{equation}
 \lambda=\frac{P\max_pW_p}{\sum_pW_p}\geq1.
\end{equation}
当工作量缩减比例为 $r$，不均衡因子却从 $\lambda_0$上升到 $\lambda_1$，简化模型下时间比例约为 $r\lambda_1/\lambda_0$。这说明选择性计算只有配合分配机制才可能稳定受益。

CV 与上述因子反映不同方面。CV 度量整体离散程度，最大值因子更接近同步查询的关键路径。两个负载集合可能 CV 相近而最大值不同。实验同时记录二者以及真实查询时间，才能区分平均分散改善与尾部拖延改善。

\subsection{预处理摊销与使用场景}
若某方案一次性准备时间为 $T_p$，单次执行时间为 $t$，连续执行 $N$ 次的总成本为 $T_p+Nt$。相对准备时间 $T_b$、执行时间 $t_b$ 的参考方案，只有
\begin{equation}
 N(t_b-t)>T_p-T_b
\end{equation}
才具有总时间优势。若 $t_b\leq t$，增加重复次数也不会弥补在线劣势。此关系适用于索引构建、数据转换和策略训练，但要求准备结果在这些执行中持续有效。

当数据库更新或矩阵变化使准备结果失效，重复次数应按两个更新事件之间的有效执行数计算，而不是系统生命周期内总请求数。尤其是 EvoCG 的训练策略可能与矩阵及求解设置相关，不能只因应用名称相同就认为训练成本永久摊销。

\subsection{综合评估的可比性条件}
不同实验应统一质量目标，却不必强行统一所有算法参数。相同探测数可能产生不同召回，相同位宽可能产生不同误差；因此，质量匹配优先于参数数值相同。为了隔离硬件差异，再补充固定算法和固定候选集合的对照。

\begin{table}[htbp]\centering\caption{跨方法比较中应固定与应报告的条件}
\begin{tabularx}{\textwidth}{p{2.5cm}XX}\toprule
比较目的&应固定条件&仍需报告的变化\\\midrule
近似算法收益&执行平台、输入、输出要求&位宽、候选预算与工作量\\
硬件架构收益&算法、候选集合、质量要求&数据布局、并行映射与接口开销\\
精度控制收益&稀疏内核与运行时基础&精度轨迹、迭代数、策略开销\\
负载均衡收益&同一查询序列与总任务&映射、峰值核负载、数据迁移\\
部署成本收益&有效使用周期与更新频率&预处理、功耗范围、复用次数\\\bottomrule
\end{tabularx}\end{table}

上述框架将质量约束、工作量变化与执行成本联系起来。本文实验覆盖其中与三项研究直接相关的条件组合，跨平台、跨任务的结果则分别在各自评估范围内解释。统一负载下的综合验证仍有待进一步开展。
